% This is text associated with the open textbook "Open Electromagnetics"
% See immediately below for author, date
% This work is licensed under the Creative Commons Attribution-ShareAlike 4.0 International License. (CC BY-SA 4.0) 
% To view a copy of the license, visit http://creativecommons.org/licenses/by-sa/4.0/

%ORB TITLE Kirchoff's Voltage Law for Electrostatics: Differential Form 
%ORB AUTHOR 0        # S.W. Ellingson, Virginia Tech (ellingson@vt.edu)
%ORB DATE 20180806   # yyyymmdd
%ORB ABSTRACT Kirchoff's Voltage Law for Electrostatics: Differential Form 

\label{m0152_KVL_Electrostatics_Integral-to-Differential} % <-- Must match name of this file and module directory name.  Don't mess with this!
\index{Kirchoff's voltage law!electrostatics|(}

The integral form of Kirchoff's Voltage Law for electrostatics (KVL; Section~\ref{m0016_Kirchoffs_Voltage_Law_Electrostatics_Integral_Form}) states that an integral of the electric field along a closed path is equal to zero:
%
\begin{equation}
\oint_{\mathcal C}{ {\bf E} \cdot d{\bf l} } = 0  
\end{equation}
%
where ${\bf E}$ is electric field intensity and ${\mathcal C}$ is the closed curve.
In this section, we derive the differential form of this equation.
In some applications, this differential equation, combined with boundary conditions imposed by structure and materials (Sections~\ref{m0020_Boundary_Conditions_on_E} and \ref{m0021_Boundary_Conditions_on_D}), can be used to solve for the electric field in arbitrarily complicated scenarios.  
A more immediate reason for considering this differential equation is that we gain a little more insight into the behavior of the electric field, disclosed at the end of this section.

The equation we seek may be obtained using Stokes' Theorem 
\index{Stokes' theorem}
(Section~\ref{m0051_Stokes_Theorem}), which in the present case may be written:
%
\begin{equation}   
\int_{\mathcal S} \left( \nabla \times {\bf E} \right) \cdot d{\bf s} = \oint_{\mathcal C} {\bf E}\cdot d{\bf l}
\label{m0152_eKVL2}
\end{equation}
%
where ${\mathcal S}$ is any surface bounded by ${\mathcal C}$, and $d{\bf s}$ is the normal to that surface with direction determined by right-hand rule.
The integral form of KVL tells us that the right hand side of the above equation is zero, so:
%
\begin{equation}   
\int_{\mathcal S} \left( \nabla \times {\bf E} \right) \cdot d{\bf s} = 0
\end{equation}
%
The above relationship must hold regardless of the specific location or shape of ${\mathcal S}$.  The only way this is possible for all possible surfaces is if the integrand is zero at every point in space.  Thus, we obtain the desired expression:
%
\begin{equation}
\boxed{
\nabla \times {\bf E} = 0
}
\label{m0152_eKVL}
\end{equation}
% 
Summarizing:
%
%ORB HIGHLIGHT
\begin{mdframed}[backgroundcolor=yellow!20,nobreak=true] 
The differential form of Kirchoff's Voltage Law for electrostatics (Equation~\ref{m0152_eKVL}) states that the curl of the electrostatic field is zero.
\end{mdframed}

Equation~\ref{m0152_eKVL} is a partial differential equation. 
As noted above, this equation, combined with the appropriate boundary conditions, can be solved for the electric field in arbitrarily-complicated scenarios.
Interestingly, it is not the only such equation available for this purpose -- Gauss' Law (Section~\ref{m0045_Gauss_Law_Differential_Form_from_Integral_Form}) also does this.
Thus, we see a \emph{system} of partial differential equations emerging, and one may correctly infer that that the electric field is not necessarily fully constrained by either equation alone.   

%-------------------------------------------------------

{\bf Additional Reading:}
\begin{itemize}
\item ``\href{https://en.wikipedia.org/wiki/Maxwell's_equations}{Maxwell's Equations}'' on Wikipedia.
\item ``\href{https://en.wikipedia.org/wiki/Boundary_value_problem}{Boundary value problem}'' on Wikipedia.
\end{itemize}

\index{Kirchoff's voltage law!electrostatics|)}

